\section{Related reviews and the gap addressed by this work}
\label{sec:related-reviews}

The intersection of surrogate modelling, optimization-based energy-
system modelling, multi-objective optimization (MOO) and multi-energy
systems (MES) sits between four parent literatures, each of which has
already produced its own review tradition. We position the present
work explicitly against those reviews. The intent is twofold: to
acknowledge that the methodological vocabulary, the application
classifications and several of the recurring observations in this
review are not novel in isolation, and to make precise where the
present work contributes a synthesis that the existing reviews leave
open. Table~\ref{tab:T7-related-reviews} compares thirty-two prior
reviews -- selected from the strict review classifier in our
bibliography (title, abstract and keyword markers, see
Appendix~\ref{app:search}) and from a curated set of externally
supplied PDFs that fall outside the Scopus search-string keyword cone
(Comp.~Chem.~Eng., Comp.~Intelligence and Neuroscience,
mathematics venues, MCDM journals; collected in
\texttt{references/external\_reviews.bib} and abstract-verified
case-by-case) -- along four scope axes (method scope, application
scope, MOO inclusion, MES inclusion) and along the auxiliary axis of
decision-aware validation that we treat as a primary diagnostic in
Section~\ref{sec:validation}. The thirty-two rows are arranged in
four blocks that mirror the four parent literatures.

% BEGIN inlined table table_T7_related_reviews
\begin{table*}[!t]
  \centering
  \scriptsize
  \caption{Meta-review of recent review studies in the surrogate $\times$
  optimization $\times$ energy-system literature, grouped by primary
  focus (surrogate methodology, multi-energy systems, multi-objective
  optimization and MCDM, energy-system optimization). ``Method
  scope'' abbreviates the surrogate or optimization families covered
  (\textsc{gp}: Gaussian process / kriging, \textsc{rsm}: response
  surface, \textsc{pce}: polynomial chaos, \textsc{rbf}: radial basis
  function, \textsc{nn}: neural network, \textsc{l2o}: learning-to-
  optimize, \textsc{piml}: physics-informed machine learning,
  \textsc{mf}: multi-fidelity, \textsc{moea}: multi-objective
  evolutionary algorithm, \textsc{mcdm}: multi-criteria decision
  making). ``MOO'' indicates whether multi-objective optimization is
  part of the review's scope; ``MES'' whether multi-energy systems
  are; ``DA-val'' whether decision-aware validation (regret,
  feasibility, OOD, calibrated tails) is treated beyond standard
  regression metrics. The last column states the gap that each
  review leaves open and that the present work therefore addresses.}
  \label{tab:T7-related-reviews}
  \renewcommand{\arraystretch}{1.15}
  \begin{tabularx}{\textwidth}{@{}L{0.10\textwidth} L{0.04\textwidth} L{0.16\textwidth} L{0.18\textwidth} C{0.04\textwidth} C{0.04\textwidth} C{0.04\textwidth} L{0.30\textwidth}@{}}
    \toprule
    Reference & Year & Method scope & Application scope & MOO & MES & DA-val & Gap left open / addressed here \\
    \midrule
    \multicolumn{8}{@{}l}{\textit{Surrogate methodology}} \\
    \cite{Bhosekar2018}
      & 2018
      & \textsc{rbf}, \textsc{gp}, derivative-free
      & Generic chemical-engineering optimization
      & part.\ & -- & part.
      & Generic methodology; no MES, no MOO Pareto search; pre-dates the L2O / decision-focused wave \\
    \cite{DiazManriquez2016}
      & 2016
      & Surrogate-assisted \textsc{moea} (8 classes)
      & Generic expensive MO benchmarks
      & yes & -- & --
      & Generic benchmarks; no energy applications; integration patterns enumerated, not formalised across host optimizers \\
    \cite{FernandezGodino2023}
      & 2023
      & \textsc{mf} surrogates (combination methods)
      & Engineering / UQ / optimization
      & part.\ & -- & part.
      & Methodology-only; energy-system applications absent; UQ-MOO interplay not isolated \\
    \cite{Westermann2019}
      & 2019
      & \textsc{gp}, \textsc{nn}, \textsc{rbf}, \textsc{rsm}
      & Sustainable building design
      & part.\ & -- & --
      & Single sector (buildings); no MES coupling; integration patterns informal \\
    \cite{Mohammadi2024}
      & 2024
      & Data-driven + \textsc{piml}
      & Optimal power flow
      & -- & -- & part.
      & OPF-only; no MOO Pareto search; no MES coverage \\
    \cite{Tan2026}
      & 2026
      & \textsc{gp} (probabilistic, UQ-focused)
      & Power-system analysis and control
      & part.\ & -- & part.
      & GP-only family; no MOO Pareto search; no MES coverage; integration patterns implicit \\
    \cite{Khan2024_DT}
      & 2024
      & AI-driven surrogate, digital twin
      & Hybrid sustainable energy systems
      & part.\ & yes & part.
      & DT framing; surrogate role broad but not split by integration pattern; MOO not isolated \\
    \cite{Khaloie2025}
      & 2025
      & ML across families
      & Optimal power flow
      & -- & -- & part.
      & Single problem class (OPF); no UC, CapEx, MES, MOO; integration patterns implicit \\
    \cite{Lim2025}
      & 2025
      & Deep learning only
      & Power system decision making
      & -- & -- & part.
      & DL-only; no PCE, GP-BO or surrogate-assisted NSGA; no MES, no MOO Pareto search \\
    \cite{Starke2025214}
      & 2025
      & ML metamodelling (general)
      & Engineering / partial energy
      & part.\ & part.\ & part.
      & General engineering scope; does not isolate ESM-specific pitfalls or MOO-MES intersection \\
    \cite{Ruan2021221}
      & 2021
      & \textsc{nn}, \textsc{gp}, RL, \textsc{l2o}
      & Power system optimization (UC, OPF, ED)
      & -- & -- & part.
      & No MES coverage; MOO not treated; integration patterns mostly enumerated, not classified \\
    \cite{Elsheikh2019622}
      & 2019
      & ANN only
      & Solar energy systems
      & -- & -- & --
      & Methodologically narrow (ANN); no MOO, no host-optimizer integration; predates surrogate-assisted MOO wave \\
    \cite{Elwy2024}
      & 2024
      & ML surrogate + \textsc{moea}
      & Sustainable building design
      & yes & -- & part.
      & Single sector (buildings); MOO present but MES absent; integration patterns not formalised \\
    \cite{Zhang2026_building}
      & 2026
      & AI surrogate (\textsc{ml}/\textsc{dl}) + \textsc{moea} + \textsc{mcdm}
      & Sustainable building design (114 studies, 5 dimensions)
      & yes & -- & part.
      & Single sector (buildings); MES coupling absent; flags DL interpretability and uncertainty handling as open gaps \\
    \midrule
    \multicolumn{8}{@{}l}{\textit{MES-domain reviews}} \\
    \cite{Mylonopoulos202332697}
      & 2023
      & Physics-based + data-driven + hybrids
      & Ship energy systems
      & part.\ & yes & part.
      & Domain-specific (ships); does not extract surrogate-MOO patterns transferable to MES at large \\
    \cite{Manco2024}
      & 2024
      & MILP / NLP / MOO formulations
      & Multi-energy systems (energy hubs, polygeneration)
      & yes & yes & --
      & Optimization-side review with detailed math formulation; surrogate role implicit \\
    \cite{Klemm2021}
      & 2021
      & MILP, LP, dynamic
      & MES in mixed-use districts (145 tools surveyed)
      & part.\ & yes & --
      & Tool inventory rather than methodology; surrogate axis not separated \\
    \cite{Fattahi2020}
      & 2020
      & Optimization + simulation suites
      & National integrated ESMs (TIMES, MARKAL, REMix, PRIMES, METIS)
      & part.\ & yes & --
      & Macro-ESM perspective; no link to surrogate-assisted intra-model acceleration \\
    \cite{agha_kassab_comprehensive_2024}
      & 2024
      & Heuristic + mathematical + hybrid
      & Microgrid sizing and energy management
      & part.\ & yes & --
      & Sizing and EMS treated jointly but at the optimization-method axis, not the surrogate axis; integration patterns not formalised \\
    \cite{nallolla_multi-objective_2023}
      & 2023
      & Evolutionary MOO algorithms
      & Hybrid AC/DC microgrid with RES
      & yes & yes & --
      & Algorithm-side MOO survey; surrogate role absent; no decision-aware validation \\
    \cite{malla_sg_optimization_2024}
      & 2024
      & MILP, MOO
      & MES with Power-to-X, P2H$_2$, district scale
      & yes & yes & --
      & Method-side MES review for P2X; surrogate use not made explicit; integration patterns not classified \\
    \cite{Li2025_hydrogen}
      & 2025
      & AI-driven + physics-assisted
      & Hydrogen-integrated HRES (solar, wind)
      & part.\ & yes & part.
      & Domain-narrow (hydrogen-HRES); MOO partial; surrogate role not abstracted into integration patterns \\
    \midrule
    \multicolumn{8}{@{}l}{\textit{MOO and MCDM methodology}} \\
    \cite{ChenRenZhou2023}
      & 2023
      & MOO methods (preference-based)
      & Long-term energy system models
      & yes & part.\ & --
      & Methods catalogue without surrogate axis; long-term horizon only \\
    \cite{salgueiro_multi-objective_2019}
      & 2019
      & MOO metaheuristics (NSGA, MOPSO, MOEA/D)
      & Microgrid planning and management
      & yes & yes & part.
      & MOO methodology focus; surrogate-assisted MOO not isolated; integration patterns informal \\
    \cite{Sahoo2025}
      & 2025
      & \textsc{mcdm} (AHP, TOPSIS, hybrid), bibliometric
      & Energy management 2010--2025
      & yes & part.\ & --
      & MCDM-only methodology axis; surrogate-assisted decision making not treated \\
    \cite{Etghani2025}
      & 2025
      & ANN + Taguchi (+ GRA, TOPSIS)
      & ICE, TES, solar/wind, heat exchangers, HRES
      & yes & part.\ & part.
      & Methods pair (ANN+Taguchi) only; integration patterns not formalised; OOD/feasibility not treated \\
    \midrule
    \multicolumn{8}{@{}l}{\textit{ESM-optimization and bibliometric reviews}} \\
    \cite{Zhou2024__2}
      & 2024
      & Sizing methods (incl.\ surrogates)
      & Renewable-storage sizing
      & part.\ & part.\ & --
      & Sizing-only; no operational MES dispatch, no chance-constrained P5 patterns \\
    \cite{vahidinasab_overview_2020}
      & 2020
      & Optimization methods (planning)
      & Distribution network expansion planning
      & part.\ & part.\ & part.
      & Planning-side review; surrogate role implicit; uncertainty handling not unified across methods \\
    \cite{Conti2026}
      & 2026
      & Forecasting + optimization (det.\ / stoch.\ / metaheuristic / hybrid)
      & Large-scale PV in smart distribution networks
      & part.\ & part.\ & part.
      & Domain-narrow (PV $\times$ distribution); MOO not principal organising axis; surrogate role not formalised \\
    \cite{arar_tahir_scientific_2023}
      & 2023
      & Bibliometric (scientific mapping)
      & Microgrids integrated with renewable energy
      & part.\ & yes & --
      & Mapping-only output; no methodological taxonomy of surrogate families and integration patterns \\
    \cite{batista_optimizing_2023}
      & 2023
      & Bibliometric + Portfolio Theory + PSO
      & Hybrid renewable energy systems for reverse osmosis
      & yes & yes & --
      & Single application (RO water); MOO with PT--PSO only; surrogate role not abstracted \\
    \cite{velasquez_intelligence_2023}
      & 2023
      & Bibliometric (clustering, co-occurrence)
      & AI techniques for sustainable energy (decade survey)
      & part.\ & part.\ & --
      & Macro-bibliometric; does not enter the surrogate $\times$ optimization-method axis at the resolution this work targets \\
    \midrule
    \multicolumn{4}{@{}l}{\textit{This work}}
      & yes & yes & yes
      & Joint coverage of all four axes; explicit five-pattern integration framework
        (Section~\ref{sec:integration}); decision-aware validation protocol
        (Section~\ref{sec:validation}); software-package landscape
        (Table~\ref{tab:T8-software-packages}) \\
    \bottomrule
  \end{tabularx}
\end{table*}
% END inlined table table_T7_related_reviews

\subsection{Surrogate-methodology reviews}
\label{sec:related-surrogate}

The surrogate-methodology line is the most mature. It traces back to
the systematic taxonomy of \cite{Bhosekar2018}, who organise surrogate
methods around three problem types -- modelling, feasibility analysis
and derivative-free optimization -- and benchmark Kriging and radial
basis functions on a battery of test problems. The same paper
contains an explicit software-implementations section that has shaped
how subsequent reviews report tooling. \cite{DiazManriquez2016}
sharpen the focus to surrogate-assisted multi-objective evolutionary
algorithms (MOEAs) and propose an integration-pattern classification
that is the conceptual ancestor of the five-pattern framework we use
in Section~\ref{sec:integration}: they distinguish surrogate-assisted
fitness approximation, model management with infill criteria, and
hierarchical evaluation, but they stay on generic MO benchmarks and
do not enter the energy domain. \cite{FernandezGodino2023} closes the
methodological frontier with a comprehensive multi-fidelity review
that classifies combination methods (linear additive, multiplicative,
hierarchical co-Kriging, deep multi-fidelity nets) and is the source
of the multi-fidelity vocabulary used throughout
Section~\ref{sec:training-doe}.

The energy-system surrogate reviews follow this methodological
backbone but each restricts the application scope. \cite{Tan2026}
covers Gaussian processes in power systems exhaustively, with a
strong focus on probabilistic prediction and uncertainty
quantification; \cite{Mohammadi2024} restricts to the optimal-power-
flow (OPF) problem and contrasts data-driven and physics-informed
approaches; \cite{Khaloie2025} reviews the broader ML-OPF field
with a similar problem-class restriction; \cite{Lim2025} surveys deep
learning across power-system decision making but excludes Gaussian
processes, polynomial chaos and surrogate-assisted NSGA;
\cite{Starke2025214} broadens to ML metamodelling of energy systems
in general but stops short of the MOO-MES intersection;
\cite{Ruan2021221} catalogues learning-assisted power-system
optimization (UC, OPF, ED) including reinforcement learning and
learning-to-optimize, and \cite{Elsheikh2019622} provides a focused
ANN review for solar systems that predates the surrogate-assisted
MOO wave. Three recent additions push the scope outward.
\cite{Khan2024_DT} place surrogate modelling inside a digital-twin
framework for hybrid sustainable energy systems and discuss the
associated uncertainty handling. \cite{Westermann2019} and
\cite{Elwy2024} document the parallel development in sustainable
building design, including 57 and 72 publications respectively that
combine surrogate models with optimization, with Latin Hypercube
Sampling and adaptive sampling on Gaussian processes as the
dominant practice. \cite{Zhang2026_building} extends this trajectory
to the most recent state of the art with a 114-study systematic
review that, uniquely among the surrogate-side reviews, organises the
field along five workflow dimensions -- parametric simulation,
dataset construction, surrogate model state of the art, optimization-
algorithm selection and \emph{multi-criteria decision-making for
design-scheme selection} -- and explicitly couples surrogate-assisted
MOO with the downstream MCDM step. The same review identifies four
open gaps that align tightly with the agenda of the present paper:
limited deep-learning exploration relative to classical ML
surrogates, lack of surrogate interpretability, insufficient handling
of uncertain parameters during optimization, and inadequate scaling
from single-building to block-level decarbonisation.

\paragraph{Where the surrogate-side reviews leave gaps.}
Three gaps are visible across all fourteen surrogate-side reviews.
First, the MOO axis is rarely organising. With the exception of
\cite{DiazManriquez2016}, \cite{Elwy2024} and \cite{Zhang2026_building},
the surrogate reviews either ignore Pareto-front search or treat it
as a closing desideratum. Second, MES coupling is absent: even
\cite{Khan2024_DT}, which mentions hybrid systems, does not separate
electricity from heat, gas or hydrogen as distinct vectors, and the
remaining reviews are restricted to a single sector. Third,
decision-aware validation -- regret, feasibility rate, calibrated
tail behaviour -- is mentioned only at the methodology level and
never turned into a structured protocol that complements the
standard regression metrics; the closest formalisation is in
\cite{Khaloie2025,Lim2025,Tan2026}, but only as ``future-work''
sections. The present work consolidates these three gaps into the
five-pattern integration framework (Section~\ref{sec:integration}),
the decision-aware validation protocol
(Section~\ref{sec:validation}), and the application synthesis
(Section~\ref{sec:applications}).

\subsection{MES-domain reviews}
\label{sec:related-mes}

The MES-domain reviews discuss optimization methods at length but
treat the surrogate as one of several modelling ingredients rather
than as an organising axis. \cite{Manco2024} is, at the time of this
writing, the most explicit attempt to give a mathematical formulation
of MES optimization that future surrogate work can attach to; the
review iterates through six case studies and identifies model
choices (MILP versus NLP, time-step granularity, thermal-network
inclusion) that change the nature of the problem and therefore
change which surrogate is admissible. \cite{Klemm2021} take a
complementary tool-inventory perspective and survey 145 modelling
tools for mixed-use districts; only thirteen are usable for MES
optimization and only two remain when open-source availability and
non-financial assessment criteria are imposed -- a finding that
motivates the software-landscape table in
Section~\ref{sec:software}. \cite{Fattahi2020} extends the inventory
to nineteen national integrated ESMs (TIMES, MARKAL, REMix, PRIMES,
METIS) and uses Multi-Criteria Analysis to identify modelling gaps
at the macro scale, but the surrogate axis is implicit (machine-
learning-assisted forecasting and load disaggregation are mentioned
only as prospective extensions). The remaining MES reviews are
domain-specific: ships in \cite{Mylonopoulos202332697}, microgrids
with PV and other renewables in \cite{agha_kassab_comprehensive_2024}
and \cite{nallolla_multi-objective_2023}, MES with Power-to-X in
\cite{malla_sg_optimization_2024}, hydrogen-integrated HRES in
\cite{Li2025_hydrogen}.

\paragraph{Where the MES-side reviews leave gaps.}
The MES reviews establish that the optimization side of the field is
mature -- MILP, NLP, MOO formulations are well-documented -- and
that surrogate-assisted variants exist but are scattered. None of
the eight MES reviews extracts a transferable surrogate $\times$
integration-pattern classification, and none links the surrogate
choice to the host optimizer. \cite{Klemm2021} is particularly clear
on this gap: of 145 tools, only two are open-source and capable of
MES optimization, but the report does not extend to which surrogate
machinery these two tools embed or whether external surrogate
libraries can be plugged in. The decision-aware validation axis is
absent in all eight MES reviews.

\subsection{MOO and MCDM-methodology reviews}
\label{sec:related-moo-mcdm}

The MOO and MCDM line is best represented by four complementary
contributions. \cite{salgueiro_multi-objective_2019} surveys multi-
objective metaheuristic challenges in microgrid planning and
management; the review is methodologically sharp on the algorithmic
side (NSGA-II/III, MOEA/D, MOPSO, hybrid evolutionary--metaheuristic
schemes) but the surrogate is not isolated as a category in its own
right. \cite{ChenRenZhou2023} broadens the lens to long-term energy-
system models and classifies MOO methods by the preferential
expression of the decision maker (a priori, interactive, a
posteriori), placing MCDM and weighted-sum methods on the same
spectrum as evolutionary search. \cite{Sahoo2025} provides the
bibliometric counterpart with a 2010--2025 review of MCDM
applications in energy management; the dominant methods are AHP,
TOPSIS and hybrid models, and the review documents the recent shift
from single-method MCDM to AI-enhanced hybrid frameworks.
\cite{Etghani2025} is the most application-driven of the four,
combining ANN surrogates with Taguchi-based design of experiments
and Grey Relational Analysis or TOPSIS for multi-objective
selection; reported gains are 20--30~\% on multi-objective scenarios
across internal-combustion engines, thermal-energy storage, solar
and wind systems, heat exchangers and HRES.

\paragraph{Where the MOO/MCDM-side reviews leave gaps.}
The MOO and MCDM reviews jointly cover the algorithm axis (NSGA,
MOEA/D, MOPSO, AHP, TOPSIS, GRA), the preference-modelling axis (a
priori, interactive, a posteriori), and the bibliometric axis. They
do not, however, isolate the surrogate as a category in its own
right. Surrogate-assisted MOO appears as one of several speed-up
strategies rather than as a class with its own integration patterns
(P1--P5) and validation needs (RMSE-to-regret, feasibility rate, OOD
robustness). The MCDM-on-Pareto-front step in particular -- where a
surrogate is interrogated for many candidate solutions before an
MCDM rule selects one -- is treated only at the workflow level and
never with the calibration questions that this paper raises in
Section~\ref{sec:validation}.

\subsection{ESM-optimization and bibliometric reviews}
\label{sec:related-esm-opt}

The optimization-side and bibliometric reviews
\cite{Zhou2024__2,vahidinasab_overview_2020,Conti2026,arar_tahir_scientific_2023,batista_optimizing_2023,velasquez_intelligence_2023}
are useful as macro-maps of the field. \cite{arar_tahir_scientific_2023}
analyses 2{,}307 Scopus records on microgrid optimization and shows
a clear shift towards MOO-flavoured studies after 2016;
\cite{velasquez_intelligence_2023} extends the lens to 18{,}715
articles on AI in sustainable energy and identifies AI-MOO and AI-
MCDM clusters as the fastest-growing topical groups;
\cite{batista_optimizing_2023} narrows to hybrid renewable systems
for reverse-osmosis water and pairs the bibliometric output with a
concrete Portfolio-Theory + PSO MOO design. \cite{Zhou2024__2}
focuses on sizing approaches for centralised and distributed
renewable storage and is the most cited recent state-of-the-art
review on the sizing axis. \cite{vahidinasab_overview_2020} and
\cite{Conti2026} cover distribution-network expansion planning and
PV integration in smart distribution networks respectively, both
flagging surrogate-assisted forecasting and chance-constrained
optimization as emerging directions.

\paragraph{Where the ESM-optimization-side reviews leave gaps.}
These reviews establish that the surrogate $\times$ MOO $\times$
MES intersection is growing fast and that it touches sizing,
capacity expansion, distribution planning and operation. They do
not, however, enter the surrogate $\times$ optimization-method
axis at the resolution that this paper targets: the bibliometric
analyses
\cite{arar_tahir_scientific_2023,batista_optimizing_2023,velasquez_intelligence_2023}
extract topical clusters and citation maps but do not produce a
methodological taxonomy of surrogate families and integration
patterns; the planning-side reviews
\cite{Zhou2024__2,vahidinasab_overview_2020,Conti2026} keep the
surrogate role implicit.

\subsection{Synthesis: three observations and the contribution gap}
\label{sec:related-synthesis}

\emph{What has been done.}
The thirty-two prior reviews jointly cover four mature research
directions: (i) surrogate methodology with a clear taxonomy of model
families, sampling strategies and accuracy diagnostics
\cite{Bhosekar2018,DiazManriquez2016,FernandezGodino2023,Westermann2019,Tan2026,Mohammadi2024,Starke2025214,Zhang2026_building};
(ii) MES optimization with detailed mathematical formulations,
tool inventories and domain studies
\cite{Manco2024,Klemm2021,Fattahi2020,agha_kassab_comprehensive_2024,nallolla_multi-objective_2023,malla_sg_optimization_2024,Li2025_hydrogen,Mylonopoulos202332697};
(iii) MOO and MCDM methodology with algorithm catalogues and
preference-modelling spectra
\cite{salgueiro_multi-objective_2019,ChenRenZhou2023,Sahoo2025,Etghani2025};
and (iv) bibliometric and macro-mapping reviews that document the
acceleration of AI-flavoured contributions
\cite{arar_tahir_scientific_2023,batista_optimizing_2023,velasquez_intelligence_2023}.
The current research focus, visible across all four blocks, is on
combining AI/ML surrogates with stochastic or robust formulations
\cite{Khaloie2025,Lim2025,Tan2026,Khan2024_DT,Li2025_hydrogen} and
on coupling them with multi-objective decision-making
\cite{DiazManriquez2016,Etghani2025,nallolla_multi-objective_2023,Sahoo2025,Elwy2024,Zhang2026_building}.

\emph{What is missing.}
\emph{First}, no prior review treats the surrogate $\times$ MOO
$\times$ MES intersection as its primary subject; the four parent
literatures each cover three out of the four scope axes at most,
and the cell that contains all four is empty. \emph{Second}, the
integration pattern that connects a surrogate to its host optimizer
is rarely the organising principle. Where the integration is
treated, it is typically split by optimization method (Bayesian
optimization, evolutionary, MILP) rather than by the role the
surrogate plays in the workflow (inner-loop replacement, outer-loop
acceleration, warm start, decomposition cut, uncertainty
propagation). The five-pattern classification in
Section~\ref{sec:integration} (P1--P5) is, to our knowledge, the
first explicit attempt to organise the surrogate-in-ESM literature
along this axis. \emph{Third}, decision-aware validation is
mentioned as a desideratum in the closing sections of the most
recent reviews
\cite{Khaloie2025,Lim2025,Starke2025214,Ruan2021221,Tan2026}
but is not turned into a structured protocol that complements the
standard regression metrics. The RMSE-to-regret argument in
Section~\ref{sec:validation} consolidates the diagnostics that the
recent literature has identified one by one and connects each
diagnostic to one of the five integration patterns.

\emph{Interesting cross-review aspects.}
Three patterns are worth highlighting. (a) The shift from black-box
ML surrogates to physics-informed and physics-guided approaches is
visible in the most recent surrogate-side reviews
\cite{Mohammadi2024,Tan2026,Khan2024_DT,Li2025_hydrogen} and is
named explicitly as the dominant near-term research direction. (b)
The open-source bottleneck is cross-cutting: \cite{Klemm2021} report
that only 2 of 145 MES tools are open-source and capable of
optimization, and \cite{Bhosekar2018} catalogue surrogate libraries
along the same dimension; we revisit this bottleneck in
Section~\ref{sec:software} and Table~\ref{tab:T8-software-packages}.
(c) Multi-fidelity is the connective tissue between the surrogate-
side and the MES-domain literatures: \cite{FernandezGodino2023}
provide the methodology, \cite{Klemm2021,Manco2024,Fattahi2020}
provide the application domains where high-fidelity simulation is
prohibitively expensive, but no review yet links the two.

The present review is therefore intended as the joint contribution
that the four parent literatures have not yet produced: a
methodological taxonomy (Section~\ref{sec:taxonomy}), an integration-
pattern classification (Section~\ref{sec:integration}), a decision-
aware validation protocol (Section~\ref{sec:validation}), an
application-level synthesis (Section~\ref{sec:applications}) and a
software-package landscape (Section~\ref{sec:software},
Table~\ref{tab:T8-software-packages}) that together bracket the
surrogate $\times$ MOO $\times$ MES intersection. The remaining
sections develop these contributions in turn.

\subsection{Software-package landscape}
\label{sec:software}

The reviewed literature consistently points at the software-package
landscape as a key bottleneck. \cite{Bhosekar2018} dedicate
Section~7 to surrogate-specific packages (DAKOTA, the SUMO
toolbox, the Surrogates Toolbox, \texttt{scikit-learn}, MATLAB's
Statistics and Machine Learning Toolbox, R-side packages, ARGONAUT)
and observe that maintenance and language coverage are uneven;
\cite{Klemm2021} flag open-source availability as a hard filter
that drops 143 of 145 MES tools and explicitly inventories
\texttt{PyPSA}, \texttt{oemof}, \texttt{Calliope},
\texttt{OSeMOSYS}, \texttt{TIMES}, \texttt{MARKAL} and
\texttt{REMix}; \cite{Mohammadi2024} note that physics-informed
surrogates depend strongly on PyTorch- or JAX-based ecosystems;
\cite{Etghani2025} report that ANN-Taguchi pipelines are typically
built in MATLAB even though the constituent methods exist in open-
source Python. Table~\ref{tab:T8-software-packages} consolidates
the resulting landscape into five blocks (general-purpose ML,
surrogate-specific, MOO and MCDM, optimization solvers, energy-
system-specific) along six attributes -- language, GPU support,
deep-learning support, maintenance status, primary use and
\emph{Sources}. The latter column is generated automatically by
\texttt{paper\_library/verify\_T8\_software\_cites.py}, which scans
the full text of every PDF in \texttt{paper\_library/fulltexts/} for
an alias of the package name and lists the cite-keys whose review
explicitly mentions the package; entries marked ``--$^\dagger$''
are those for which no review-side mention was found in our PDF
corpus and that we kept on common-knowledge grounds (typical
example: \texttt{pymoo}, \texttt{Pyomo}, \texttt{JuMP} -- documented
in their own publications and widely used in the underlying
primary literature, but not named in any of the surveyed reviews).
The companion CSV \texttt{paper\_library/software\_landscape.csv}
records the same data in machine-readable form, including the
``confidence'' flag (``high'' if at least one PDF mention, ``ck''
otherwise) used to populate the table.

% BEGIN inlined table table_T8_software_packages
\begin{table*}[!t]
  \centering
  \scriptsize
  \caption{Software-package landscape for surrogate modelling, multi-objective and multi-criteria optimization, and energy-system optimization. \textit{Sources} lists the reviews from Table~\ref{tab:T7-related-reviews} whose full text explicitly mentions the package; the mapping is generated automatically from \texttt{paper\_library/verify\_T8\_software\_cites.py}, which scans the corresponding PDFs in \texttt{paper\_library/fulltexts/} for an alias of the package name. Entries marked ``--$^\dagger$'' are kept on common-knowledge grounds (the package is documented in the package's own publication or repository but is not named in any of the surveyed reviews); they are flagged here so the reader can distinguish review-documented practice from broader engineering practice. ``GPU'' indicates whether GPU acceleration is a first-class feature (yes), present via a dependency (part.), or absent (--). ``DL'' indicates whether deep-learning models are first-class (yes), present but not the focus (part.), or absent (--). ``Maint.''\ is the maintenance status as of 2026 (active, semi-active, dormant, commercial), drawn from the package's public repository. All packages are open-source unless marked as commercial. License and version pins are kept in \texttt{paper\_library/software\_landscape.csv}.}
  \label{tab:T8-software-packages}
  \renewcommand{\arraystretch}{1.10}
  \begin{tabularx}{\textwidth}{@{}L{0.13\textwidth} L{0.06\textwidth} C{0.04\textwidth} C{0.04\textwidth} L{0.07\textwidth} L{0.42\textwidth} L{0.18\textwidth}@{}}
    \toprule
    Package & Language & GPU & DL & Maint. & Primary use & Sources \\
    \midrule
    \multicolumn{7}{@{}l}{\textit{General-purpose ML / surrogate frameworks}} \\
    \texttt{scikit-learn} & Python & -- & -- & active & Linear/kernel/tree regressors used as cheap surrogates & \cite{Tan2026,Bhosekar2018,Westermann2019,Khan2024_DT} \\
    \texttt{XGBoost} & Python, R & yes & -- & active & Gradient-boosted ensembles for forecasting and infill modelling & \cite{Khaloie2025,Conti2026,Zhang2026_building} \\
    \texttt{LightGBM} & Python & yes & -- & active & Fast gradient-boosted ensembles for high-dim ESM features & \cite{Lim2025,Elwy2024} \\
    \texttt{TensorFlow} & Python, C++ & yes & yes & active & Generic NN surrogate stack & \cite{Tan2026,FernandezGodino2023} \\
    \texttt{PyTorch} & Python & yes & yes & active & NN, GP-on-PyTorch and PINN backbones & \cite{Lim2025,Westermann2019,FernandezGodino2023} \\
    \texttt{Keras} & Python (TF) & yes & yes & active & High-level API for ANN surrogates in building / HRES studies & --\textsuperscript{$\dagger$} \\
    \texttt{JAX} & Python & yes & yes & active & Differentiable programming for PINN and large-scale BO & --\textsuperscript{$\dagger$} \\
    \multicolumn{7}{@{}l}{\textit{Surrogate-specific libraries (GP, PCE, RBF, MF, BO, UQ)}} \\
    \texttt{SMT} & Python & -- & -- & active & GP, RBF, RMTS, mixture-of-experts; multi-fidelity & \cite{Westermann2019,Khan2024_DT} \\
    \texttt{GPy} & Python & -- & -- & semi-active & Vanilla Gaussian-process regression and classification & \cite{Tan2026} \\
    \texttt{GPyTorch} & Python (PyTorch) & yes & part. & active & Scalable / variational GPs on GPU & \cite{Tan2026} \\
    \texttt{botorch} & Python (PyTorch) & yes & part. & active & Bayesian optimization, qEHVI for MOO Pareto search & \cite{Tan2026} \\
    \texttt{chaospy} & Python & -- & -- & active & Polynomial chaos expansion and Sobol indices & --\textsuperscript{$\dagger$} \\
    \texttt{OpenTURNS} & Python, C++ & -- & -- & active & UQ / PCE / Kriging / sensitivity & --\textsuperscript{$\dagger$} \\
    \texttt{Modulus} & Python & yes & yes & active & NVIDIA framework for PINN and physics-guided NN & --\textsuperscript{$\dagger$} \\
    \texttt{DeepXDE} & Python & yes & yes & active & PINN library for forward and inverse PDE problems & --\textsuperscript{$\dagger$} \\
    \multicolumn{7}{@{}l}{\textit{Multi-objective optimization and MCDM}} \\
    \texttt{pymoo} & Python & -- & -- & active & NSGA-II/III, MOEA/D, MOPSO, RVEA, indicator metrics & --\textsuperscript{$\dagger$} \\
    \texttt{DEAP} & Python & -- & -- & semi-active & Evolutionary algorithm framework with multi-objective support & --\textsuperscript{$\dagger$} \\
    \texttt{pyDecision} & Python & -- & -- & active & MCDM (AHP, TOPSIS, VIKOR, PROMETHEE, ELECTRE) & --\textsuperscript{$\dagger$} \\
    \multicolumn{7}{@{}l}{\textit{Optimization solvers and modelling languages}} \\
    \texttt{Pyomo} & Python & part. & -- & active & Algebraic modelling for MILP/NLP/MINLP & --\textsuperscript{$\dagger$} \\
    \texttt{JuMP} & Julia & part. & -- & active & Julia algebraic modelling for MILP/NLP/MINLP & --\textsuperscript{$\dagger$} \\
    \texttt{CVXPY} & Python & part. & -- & active & Convex optimization for surrogate-feasibility projection & --\textsuperscript{$\dagger$} \\
    \texttt{GAMS} & GAMS & -- & -- & commercial & Algebraic modelling backbone for TIMES, MARKAL, REMix & \cite{agha_kassab_comprehensive_2024,malla_sg_optimization_2024,Fattahi2020,Klemm2021,ChenRenZhou2023} \\
    \texttt{Gurobi} & solver & -- & -- & commercial & Backend MILP/QP/SOCP solver & \cite{Lim2025,Manco2024} \\
    \texttt{CPLEX} & solver & -- & -- & commercial & Backend MILP/QP/SOCP solver & \cite{Lim2025,arar_tahir_scientific_2023,Manco2024,Fattahi2020} \\
    \texttt{Mosek} & solver & -- & -- & commercial & Backend QP/SOCP/conic solver & \cite{Mylonopoulos202332697,Fattahi2020} \\
    \texttt{IPOPT} & solver & -- & -- & active & Open-source NLP solver behind surrogate-feasible glue code & \cite{Khaloie2025,Westermann2019} \\
    \texttt{BARON} & solver & -- & -- & commercial & Global solver for non-convex MINLP and surrogate-DFO loops & \cite{Bhosekar2018} \\
    \multicolumn{7}{@{}l}{\textit{Energy-system-specific frameworks}} \\
    \texttt{PyPSA} & Python & -- & -- & active & Power-system optimization (UC, OPF, CapEx) with stochastic extensions & \cite{Klemm2021} \\
    \texttt{oemof} & Python & -- & -- & active & Modular MES framework (oemof.solph) for sector coupling & \cite{Klemm2021} \\
    \texttt{Calliope} & Python & -- & -- & active & Multi-scale MES optimization with renewables and storage & \cite{Klemm2021} \\
    \texttt{OSeMOSYS} & GAMS, Python & -- & -- & active & Long-term capacity expansion and operations & \cite{Fattahi2020,Klemm2021} \\
    \texttt{HOMER Pro} & proprietary & -- & -- & commercial & Microgrid sizing and HRES techno-economic optimization & \cite{agha_kassab_comprehensive_2024,Mylonopoulos202332697,nallolla_multi-objective_2023,velasquez_intelligence_2023,arar_tahir_scientific_2023,batista_optimizing_2023,Klemm2021,Li2025_hydrogen,ChenRenZhou2023} \\
    \texttt{TIMES} & GAMS & -- & -- & active & Long-term integrated ESM (national / EU scale) & \cite{Fattahi2020,Klemm2021} \\
    \texttt{MARKAL} & GAMS & -- & -- & active & Long-term integrated ESM (precursor of TIMES) & \cite{Fattahi2020,Klemm2021,ChenRenZhou2023} \\
    \texttt{REMix} & Python, GAMS & -- & -- & active & Integrated ESM with high temporal/spatial resolution & \cite{Fattahi2020,Klemm2021} \\
    \texttt{EnergyPlus} & C++ & -- & -- & active & Whole-building energy simulation; reference simulator & \cite{Westermann2019,Klemm2021,Zhang2026_building} \\
    \texttt{TRNSYS} & proprietary & -- & -- & commercial & Transient simulation of building / HVAC / renewable systems & \cite{Elsheikh2019622,batista_optimizing_2023,Westermann2019,Klemm2021,Li2025_hydrogen,Etghani2025} \\
    \texttt{Modelica/Dymola} & Modelica & -- & -- & active & Multi-physics simulation for thermal MES & \cite{Klemm2021,Elwy2024} \\
    \texttt{MATLAB} & MATLAB & part. & part. & commercial & General host environment for ANN, Taguchi and DoE pipelines & \cite{Tan2026,agha_kassab_comprehensive_2024,Mylonopoulos202332697,nallolla_multi-objective_2023,Lim2025,arar_tahir_scientific_2023,Bhosekar2018,Westermann2019,Li2025_hydrogen,Khan2024_DT,FernandezGodino2023,Elwy2024,Zhang2026_building} \\
    \bottomrule
  \end{tabularx}

  \vspace{2pt}
  \footnotesize
  $^\dagger$ No mention of the package was found in the full text of
  the thirty-one review PDFs in \texttt{paper\_library/fulltexts/}.
  The entry is retained because the package is part of the de-facto
  practice for surrogate $\times$ MOO $\times$ MES workflows
  (e.g.\ \texttt{pymoo}, \texttt{Pyomo}, \texttt{JuMP}) and
  removing it would obscure the open-source bottleneck reported in
  \cite{Klemm2021,Bhosekar2018}.
\end{table*}
% END inlined table table_T8_software_packages

The table reveals three structural patterns. \emph{First}, the
surrogate-specific Python ecosystem (\texttt{SMT}, \texttt{GPy},
\texttt{GPyTorch}, \texttt{botorch}, \texttt{chaospy},
\texttt{OpenTURNS}, \texttt{Modulus}, \texttt{DeepXDE}) is
fragmented: no single package covers the GP--PCE--RBF--PINN
continuum, and the GPU-aware GP fork (\texttt{GPyTorch},
\texttt{botorch}) lives inside the deep-learning ecosystem rather
than the classical surrogate ecosystem. \emph{Second}, the energy-
system tools (\texttt{PyPSA}, \texttt{oemof}, \texttt{Calliope},
\texttt{OSeMOSYS}, \texttt{TIMES}, \texttt{MARKAL}, \texttt{REMix})
almost universally lack a built-in surrogate interface; the
surrogate has to be wired into the host optimizer through ad-hoc
Python or Julia glue code. \emph{Third}, the review-mention
distribution is itself revealing: \texttt{HOMER}, \texttt{MATLAB}
and \texttt{GAMS} dominate (9, 13 and 5 review mentions
respectively), while the entire MOO and MCDM Python block
(\texttt{pymoo}, \texttt{DEAP}, \texttt{pyDecision}) carries no
review mention at all -- a quantitative confirmation of the
``algorithm-axis-without-software-axis'' gap discussed in
Section~\ref{sec:related-moo-mcdm}. All three observations
reinforce the open-challenge agenda in
Section~\ref{sec:open-challenges} (open-source tooling for
decision-aware surrogates, public datasets, MES-aware MOO
benchmarks).
